% 本文件由 LaTeX 排版 Agent 根据有效 translation.json 自动生成。
% author=Councils; work_id=3812; title=第二次君士坦丁堡公会议

\chapter{第二次君士坦丁堡公会议}

% structure: p0001 source=h1 target=omitted reason=page-title-already-rendered-by-parent

% structure: p0002 source=h2 target=section reason=relative-heading-rank; base=h2

\section{第一届会议——会议记录摘录}

[\Emph{在教父们面前宣读的皇帝信函}]\par

因我们的主天主耶稣基督之名。皇帝弗拉维·查士丁尼，日耳曼的、哥特的等等，永为奥古斯都，致至福的主教和宗主教们：\Place{君士坦丁堡}的\Person{尤提基乌斯}，\Place{亚历山大}的\Person{阿波利纳里乌斯}，\Place{特奥波利斯}的\Person{多米尼努斯}，\Person{斯德望}、\Person{乔治}、\Person{达米安}，这些代表那位独一无二的至福之人、\Place{耶路撒冷}总主教兼宗主教\Person{尤斯托基乌斯}的至虔诚的主教，以及来自不同省份停留在这座皇城中的其他至虔诚的主教们。\par

[\Emph{以下是信函的缩略版本，包括\Person{黑菲勒}的摘要}。\WorkTitle{大公会议史}，\Emph{第IV卷，第298页}]\par

我的前任们，即信奉正统信仰的皇帝们，始终致力于通过召开主教会议来解决关于信仰而起的争议。为此，君士坦丁在尼西亚召集了三百一十八位教父，并亲自出席该大公会议，协助那些宣认圣子与圣父同一质的人。狄奥多西在君士坦丁堡召集了一百五十位教父，小狄奥多西召开了厄弗所大公会议，马尔基安皇帝则在加采东召集了主教们。然而，在马尔基安去世后，关于加采东大公会议的争议在多个地方爆发，皇帝利奥致信各地所有主教，要求每个人就这一神圣大公会议书面陈明自己的意见。但不久之后，内斯多略和欧提克的追随者再次兴起，造成了巨大的分裂，以致许多教会彼此断绝了共融。如今，天主的恩宠将我们提拔至皇位，我们将重新联合各教会、并使加采东大公会议与前三次大公会议一同为普世接受视为首要事务。我们赢得了许多先前反对该大公会议的人；其他坚持反对者，我们予以放逐，从而再次恢复了教会的合一。然而，内斯多略派企图将他们的异端强加于教会；由于他们不能为此目的利用内斯多略本人，便急忙通过内斯多略的老师、摩普绥的狄奥多来引入他们的错误，此人教导的亵渎言论比他老师更严重。例如，他坚持说天主圣言是一回事，基督是另一回事。同样地，他们利用了狄奥多莱那些不敬的著作，这些著作是针对第一次厄弗所大公会议、针对济利禄及其十二章的，还有所谓依巴斯写的可耻信件。他们声称这封信被加采东大公会议接受，从而企图使内斯多略和狄奥多免受谴责，因为信中称赞了他们。倘若他们得逞，就不能再说圣言\Quote{成了人}，也不能称玛利亚为天主之母（\OriginalTerm{母亲}{genetrix}）。因此，我们追随圣教父们，首先以书面请求你们对所称的三个不敬章节作出判断，你们已经回应并怀着喜悦宣认了真信仰。然而，在你们作出谴责之后，仍有人为这三章案辩护，故我们召你们来到京城，以便你们在此共同集会，再次将你们的观点公之于众。例如，当古罗马教宗维吉里来到此处时，他答复我们的询问，多次以书面绝罚了三章案，并通过许多方式证实了他对此观点的坚定，甚至谴责了他的执事路斯提古和塞巴斯蒂安。我们仍持有他亲笔的声明。随后他发布了《判决书》，在其中绝罚了三章案，且使用了\Quote{Et quoniam}等词语。你们知道，他不仅因路斯提古和塞巴斯蒂安为三章案辩护而将他们撤职，还致信斯库提亚主教瓦伦提尼安和阿尔勒主教奥勒良，要求不得采取任何反对《判决书》的行动。当你们后来应我的邀请来到这里时，你们与维吉里之间交换了信函，以求共同集会。但现在他改变了观点，不再愿意召开主教会议，而只要求三位宗主教和另一位主教（与教宗共融及其身边的三位主教）来决定此事。我们多次命令他参加主教会议，均属徒劳。他也拒绝了我们提出的两个建议：要么设立一个法庭来裁决，要么举行较小规模的会议，会上除他和他的三位主教外，每位宗主教及其教区的三至五位主教均应有席位和发言权。* 我们进一步声明，我们坚持四项大公会议的法令，并在一切事上追随圣教父们：亚大纳削、依拉利、巴西略、额我略·纳齐盎、尼撒的额我略、盎博罗削、德奥斐洛、君士坦丁堡的若望（金口）、济利禄、奥斯定、普洛科鲁斯、良，以及他们论真信仰的著作。然而，由于异端者决心为摩普绥的狄奥多和内斯多略及其不敬行为辩护，并声称依巴斯的信被加采东大公会议接受，所以我们劝勉你们关注狄奥多的不敬著作，尤其是他在厄弗所和加采东被提出、并被两次大公会议连同那些如此持守或仍然持守的人一同绝罚的《犹太信经》；我们还劝勉你们思考圣教父们关于他和他的亵渎言论所写的，以及我们前任所颁布的，还有教会史家关于他所陈述的。你们由此将看到，他和他的异端早已被谴责，因此他的名字久已从摩普绥教会的双联礼典中被删除。请考虑那种荒谬的主张，即异端者死后不应被绝罚；我们进一步劝勉你们在此事上追随圣教父们的教导，他们不仅谴责活着的异端者，而且也在死后绝罚那些死于不义的人，正如他们为那些不公正被判罪的人死后平反、并将他们的名字写入神圣双联礼典一样；这曾发生在对君士坦丁堡主教、虔诚纪念的若望和弗拉维亚诺身上。此外，我们劝勉你们审查狄奥多莱的著作和据称依巴斯的信，其中否认圣言的降生成人，拒绝\Quote{天主之母}这一称谓和神圣的厄弗所大公会议，称济利禄为异端者，并辩护和赞扬狄奥多与内斯多略。而由于他们说加采东大公会议接受了这封信，你们必须将该大公会议有关信仰的声明与这封不敬信的内容加以比较。最后，我们恳请加速此事。因为，当被问及正确信仰时，若迟迟不作答，无异于否认正确信仰。因为在信仰之事上的提问与回答，不是谁先谁后的问题，而是谁的宣认更正、更愿悦乐天主的问题。愿天主长久保存你们，至圣至善的教父们。于五月二日，在君士坦丁堡，至尊皇帝犹斯提尼安、永久的奥古斯都登基第二十七年，最显赫的巴西略执政后第十二年。\par

% structure: p0007 source=h2 target=section reason=relative-heading-rank; base=h2

\section{第七次会议——摘自《会议记录》}

\EditorialNote{（来自巴黎手稿，见于Hardouin\WorkTitle{大公会议集}，卷三，171页以下；Mansi，卷九，346页以下。此篇演说未在其他任何手稿中完整保存。巴拉林尼（Hefele注）对增补部分的真实性提出异议[见于Noris\WorkTitle{著作集}，卷四，1037页]，但Hefele认为这些异议无关紧要。[\WorkTitle{大公会议史}，卷四，323页，注2。]所有手稿一致记载：最荣耀的圣宫财务官君士坦丁受最虔诚的皇帝派遣，他进入会议后发言如下：\Quote{Certum est vestræ beatitudini, quantum, etc.}其余部分在各手稿中互有差异。我依照巴黎本。]）}\par

你们知道，战无不胜的皇帝始终何等关切因某些人对\OriginalTerm{三章}{Three Chapters}所起的争论应有一个终结……为此，他要求最虔诚的维吉略与你们一同集会，并按照正统信仰就此事拟订一项法令。因此，虽然维吉略已多次以书面形式谴责\OriginalTerm{三章}{Three Chapters}，且曾在皇帝、最荣耀的法官及本会议许多成员面前亲口如此行，并始终准备对摩普绥提亚的狄奥多尔的维护者、号称依巴斯所写的信函，以及狄奥多勒所写反对正统信仰和圣济利禄\OriginalTerm{十二条款}{twelve capitula}的著作予以\OriginalTerm{绝罚}{anathema}，然而他却拒绝在与你及你的会议共融中做此事。\par

昨日，维吉略派遣罗马教会最可敬的副辅祭塞尔乌斯·德伊，邀请贝利撒留、策泰古斯，以及最荣耀的执政官尤斯蒂努斯和君士坦丁，还有主教狄奥多尔、阿西达斯、贝尼格努斯和福卡斯到他那里去，因为他想通过他们向皇帝传达一个答复。他们去了，但很快返回并告知最虔诚的君主：\Quote{我们拜访了最虔诚的主教维吉略，他对我们说：\InnerQuote{我们召唤你们来，是为使你们知道过去几天所行之事。为此，我就有争议的\OriginalTerm{三章}{Three Chapters}写了一份文书给最虔诚的皇帝，请你们阅读并呈送陛下。}然而当我们听闻此事并见到这份写给陛下的文书时，我们对他说：\InnerQuote{没有陛下的命令，我们绝不能接收任何写给最虔诚的皇帝的文书。但你们有跑腿的辅祭，可以藉他们发送。}他则对我们说：\InnerQuote{你们现在知道我制作了这份文书。}我们主教们回答他说：\InnerQuote{若蒙尊驾愿意与我们及至圣宗主教们和最虔诚的主教们会合，共同商讨三章问题，并与我们全体一致地制订正统信仰的适当格式，正如圣宗徒、圣教父和四大公会议所行的那样，我们愿以你为首、为父、为首席。但若尊驾已为皇帝拟订了文书，如我们所说，你们有信差，可交由他们送去。}他听了我们这些话后，就派出了副辅祭塞尔乌斯·德伊，他如今正等候陛下的答复。}陛下听闻此事后，便通过前述最虔诚、最荣耀的人士命令那位副辅祭给最虔诚的维吉略传回此言：\Quote{我们邀请他（你）与至福的宗主教们及其他虔诚的主教们会合，共同审查并判断三章问题。但你既拒绝如此行，并声称你独自就三章问题写了一些东西；若你按照你先前所行的谴责了它们，我们已有很多此类声明，不再需要；但若你现在所写与你之前所为相悖，则你以你自己的文字定自己的罪，因为你已背离正统道理，而维护了不虔。你怎能期望我们从你那里接收这样的文书呢？}\par

最虔诚的皇帝作出此答复后，他没有通过那位辅祭送来任何自己的书面文书。这一切都是在没有书面记录的情况下行事的，同样也向尊驾通报。\par

\EditorialNote{（接着，他按照所有手稿提交了一些文件，以便宣读；在Labbe和Cossart印刷的手稿中，卷五，第549栏以下。这些文件少于巴黎本，而巴黎本在宣读文件之后、会议宣布文件证明了皇帝对信仰的热忱之后还包含以下内容。）}\par

最荣耀的财务官君士坦丁说：我仍因向你们宣读的文件而出席你们的圣会议，我要说：最虔诚的皇帝向你们的圣会议发出了一份\OriginalTerm{照会}{formam}，涉及维吉略的名字，即不可再将其载入教会的圣\OriginalTerm{礼书}{diptychs}，因为他的不虔。你们不可在圣礼中提及它，也不可在都城教会或其他托付给你们及受天主托付给他治理的国家中其他主教的教会中保留它。当你们听到这份照会时，你们将再次从中感知，最温和的皇帝多么关心圣教会的合一及圣奥迹的纯洁。\par

\EditorialNote{（接着宣读信件。）}\par

神圣的会议说：最虔诚的皇帝所定之事，与他为教会合一所付的辛劳相符。让我们保持与古罗马最神圣教会宗座的合一，按照所宣读的内容完成一切。至于\OriginalTerm{所提议的问题}{De proposita vero quaestione quod jam promisimus procedat}，我们已承诺的，将继续进行。\par

% structure: p0016 source=h2 target=section reason=relative-heading-rank; base=h2

\section{会议的判决}

我们的伟大天主和救主耶稣基督，正如我们从福音中的比喻所学到的，按各人的能力分给每人塔冷通，并在适当的时候要求每人交账。如果那领了一个塔冷通的人，因没有用它工作而只是保持无损失就被定罪，那么，那不仅自己疏忽，而且在他人面前放置绊脚石和使人跌倒之缘由的人，将受到何等更大更可怕的审判呢？因为所有信徒都明瞭，每当信仰方面发生任何问题时，不仅不敬神的人自己被定罪，而且那有权纠正他人不敬神却忽略这样做的人，也要被定罪。\par

因此，我们受命治理主的教会，害怕那加于疏忽主工的人身上的诅咒，急忙保持信仰的好种子纯洁，不受敌人所播种的不敬神稗子的污染。\par

因此，当我们看到聂斯多略的追随者试图通过不敬神的\Person{戴奥多若}（他曾是\Place{摩普绥提亚}的主教）及其不敬神的著作，以及通过\Person{狄奥多勒}不敬神地所写的那些东西，和据说由\Person{依巴斯}写给波斯人\Person{马里斯}的邪恶书信，将他们的不敬神引入天主的教会时，受所有这些景象的推动，我们起来纠正所发生的事，并在这座皇城聚集，这是应天主的旨意和最虔诚皇帝的谕令而来的。\par

因为最虔诚的\Person{维吉流}恰好停留在这座皇城，亲自参与了一切关于三章案的讨论，并曾多次口头和书面谴责它们，然而后来他书面同意出席大公会议并与我们一同审查三章案，以便由我们所有人提出一个关于正确信仰的合适定义。此外，最虔诚的皇帝，按照我们之间所商定的，劝勉他和我们会面，因为司铎职在共同讨论后制定共同信仰是合宜的。为此，我们恳请他的尊贵履行他的书面承诺；因为关于这三章案的丑闻不应再继续下去，天主的教会也不应因此受到扰乱。为此，我们使他想起宗徒们留给我们的伟大榜样和教父们的传统。因为虽然圣神的恩宠充沛于每位宗徒身上，以致他们中无人需要在执行工作时借助他人的忠告，但他们仍不愿在当时提出的关于外邦人割礼的问题上作出决定，直到他们聚集在一起，用圣经的见证确认了他们各自的说法。\par

于是他们一致达成这句判决，写给外邦人说：\BibleQuote{圣神和我们决定，不再将别的重担加给你们，唯有这几件事是必要的：就是禁食祭偶像之物、血、勒死的牲畜和奸淫。}\par

而且，历代在四次神圣大公会议中聚会的圣教父们，仿效古人的榜样，通过共同讨论，以确定的法令处置了已出现的异端和问题，因为确知，当双方提出争议问题时，通过共同讨论，真理之光驱散谎言之黑暗。\par

此外，在有关信仰的讨论中，没有其他方式能使真理显明，因为每人需要邻人的帮助，正如我们在撒罗满的\BibleBook{}{箴言}中所读到的：\BibleQuote{兄弟帮助兄弟，必如设防的城，坚固如稳固的王国}；又在\BibleBook{}{训道篇}中他说：\BibleQuote{二人胜于一人，因为他们劳作的报酬美好。}\par

主自己也这样说：\BibleQuote{我实在告诉你们：若你们中二人在地上同心合意，无论为任何事祈求，我在天之父必成全他们。因为无论在哪里，有两三个人因我的名聚集，我就在他们中间。}\par

然而，当我们众人屡次邀请他，最光荣的法官也被最虔诚的皇帝派遣给他时，他承诺亲自就三章案\OriginalTerm{发出判决}{sententiam proferre}。听到这个答复后，我们想起宗徒的劝告，即\BibleQuote{每人必须向天主交自己的账}，并害怕那加于使最小者之一跌倒之人的审判，且知道冒犯如此完全基督徒的皇帝、民众和所有教会是多么严重；此外，也回忆天主对保禄所说的话：\BibleQuote{不要怕，只管讲，不要缄默，因为我与你同在，没有人能伤害你。}因此，聚集在一起后，我们先简略宣认：我们持守那由我们的主耶稣基督、真天主，传给祂的圣宗徒，并通过他们传给圣教会的信仰，以及此后继承他们的圣教父和圣师们交托给他们所照顾的民众的信仰。\par

我们宣认，我们持守、保存并向圣教会宣告那由聚集在\Place{尼西亚}的三百一十八位圣教父更详尽阐述的信德宣认，他们传授了圣 \OriginalTerm{信经}{mathema} 或信经。此外，聚集在\Place{君士坦丁堡}的一百五十位教父阐明了我们的信德，他们遵循并解释了同一信德宣认。而在第一次\Place{厄弗所}大公会议中，为同一信德而聚集的二百位圣教父也表示赞同。还有在\Place{加采东}为同一信德而聚集的六百三十位教父所定义的内容，他们遵循并教导了同一信德。所有曾由天主教会及上述四次大公会议谴责或绝罚的人，我们宣认他们被谴责和绝罚。当我们如此宣认信德后，便开始审查《三章》之事。首先，我们审查了\Person{莫普苏埃斯提亚的狄奥多若}的问题；当他的著作中所含的一切亵渎之词被揭露时，我们惊叹于天主的容忍，那编造如此亵渎的舌头和心智竟未立即被神火烧毁；我们绝不愿任由上述亵渎之词的读者继续下去——我们惧怕天主因记录这些亵渎而发怒（因为每一亵渎在不敬的程度上都超越前一个，动摇了听者的根基）——若非我们看到那些以此类亵渎为荣的人正需要因这些亵渎被揭露而遭受的羞愧。于是，我们全体在阅读期间及之后，因这些亵渎天主之词而义愤填膺，纷纷对狄奥多若发出谴责和绝罚，仿佛他还活着且在眼前。我们喊道：\Quote{主啊，求你仁慈，连魔鬼也不敢对你说出这样的话。}\par

哦，无法容忍的舌头！哦，此人的败坏！哦，他高举手臂攻击他的造物主！因为那可怜人曾许诺认识圣经，却不记得\Person{欧瑟亚}先知的话：\BibleQuote{祸哉，因为他们从我逃遁！他们因对我行恶而成名；他们说恶言反对我，他们设想后，说了强暴的话反对我。因此他们必因自己舌头的邪恶而陷于罗网。他们的轻慢必归到自己怀中，因为他们违背了我的盟约，并对我法律行事不虔敬。}\par

不虔敬的狄奥多若正应受这些咒骂。因为他拒绝关于基督的预言，并竭力摧毁我们救恩的\OriginalTerm{救恩计划的经济}{dispensation}的伟大奥迹，想方设法证明圣言不过是寓言，为外邦人取笑；又蔑视其他针对不虔敬之人所作的预言，尤其是神圣的\Person{哈巴谷}论及那些错误教导的人时所说的：\BibleQuote{祸哉，那给予邻舍饮酒，并向托住他的器皿倒酒，使他陶醉，好窥看他们的裸体的人}，即他们充满黑暗、完全与光明隔绝的教义。\par

我们为何还要多加什么？因为任何人都可以拿起不虔敬的狄奥多若的著作，或我们从他不虔敬的著作中插入我们法案中的不虔敬章节，发现他所说的难以置信的愚蠢和可憎之事。因为我们惧怕进一步继续，又记起这些丑行。\par

还向我们宣读了圣教父们反对他的著作，以及他那超越所有异端者的愚昧，此外还有历史和帝国法律，从一开始就表明他的不虔敬。既然在这一切之后，他邪恶的辩护者们以他对造物主所说的伤害为荣，说死后不应绝罚他——尽管我们知道关于不虔敬者的教会传统，即异端者即使死后也受绝罚——我们仍认为有必要对此加以审查。在法案中发现，许多异端者曾被死后绝罚。我们清楚地看到，那些说这话的人毫不关心天主的审判，也不关心宗徒的宣告和教父的传统。我们想问他们，对主关于自己所说的话，他们有何可说：\BibleQuote{凡信他的人，不被审判；那不信的人，已被审判了，因为他没有信天主子独生子的名字}，以及宗徒的宣告：\BibleQuote{即使是我们，或是从天上来的天使，若给你们宣讲我们所传给你们的福音以外的福音，当受绝罚。正如我们先前说过的，现在我再说一遍：若有人给你们宣讲你们所领受的福音以外的福音，当受绝罚。}\par

因为当主说\BibleQuote{他已受审判}，当宗徒绝罚天使（若他们教导我们所宣讲以外的任何东西），那些胆大妄为的人怎能声称这些话只针对活人？还是他们不知道，或不如说他们装作不知道，绝罚的判决无非是与天主分离？因为不虔敬的人，即使无人以言语绝罚他，他实际上已被绝罚，因他的不虔敬使自己与真生命分离。\par

再者，对于宗徒说的：\BibleQuote{异端者，在第一次和第二次劝戒后，当弃绝。因为你知道这样的人已败坏，且犯罪，而为自己定了罪}，他们有何可说？\par

根据这些话，可敬的\Person{济利禄}于其驳斥\Person{狄奥多尔}的著作中说了如下的话：\Quote{应当躲避那些陷入如此可怕罪行的人，无论他们是否在世。因为必须始终避开有害之事，而不应看重人的情面，只应考虑什么是中悦天主的。}同样，同一位可敬的\Person{济利禄}，写信给\Place{安提约基雅}主教\Person{若望}及在该城集会的教务会议，论及与\Person{奈斯多利}一同受绝罚的\Person{狄奥多尔}，如此说：\Quote{因此必须举行辉煌的庆典，因为所有附和\Person{奈斯多利}亵渎言论的言论都已被驱逐，无论出自何人。因为它针对所有持有同样见解或曾经持有的人，这正合我们与你的圣德所说：\InnerQuote{我们绝罚那些说有两个儿子、两个基督的人。}因为被我们和你们所宣讲的，正如我们所说，是一位，基督、子、主，作为人的独生子，按照博学至极的\Person{保禄}的话。}另外，在他给\Person{亚历山大}、\Person{玛尔提尼亚诺}、\Person{若望}、\Person{帕雷戈里}和\Person{玛克西姆}（司铎、隐修院长及与他们一同度独修生活者）的信中，他如此说：\Quote{\Place{厄弗所}大公会议，按天主的旨意聚集起来，反对\Person{奈斯多利}的背信，以公正而严厉的判决，与他一同谴责了那些日后支持或昔日曾支持他同样意见的人，并胆敢说或写任何此类事物，对他们施以同等的谴责，连同他们的虚妄言论。因为自然的结果是，当一人因如此亵渎的虚妄言论而被定罪时，判决不应仅针对一人，而应（可说是）针对他们整个的异端或诽谤，即他们对基督虔诚教义所发出的，崇拜两个儿子，分裂不可分者，并将人崇拜（\OriginalTerm{人崇拜}{anthropolatry}）的罪行带入天上和地上。因为同我们一道，天上诸圣的大众钦崇一位主耶稣基督。}此外，还宣读了几封信件，出自至敬礼的\Person{奥斯定}，他在\Place{非洲}主教中熠熠生辉，这些信件表明，异端者死后受绝罚是极其正确的。而这一教会传统，\Place{非洲}其他最可敬的主教们一直保存着；并且圣\Place{罗马}教会也曾绝罚过某些主教死后，尽管他们生前未曾被控有任何背信行为。关于每一条，我们手头都有证据。\par

然而，由于特奥多雷及其不敬之心的门徒——他们显然是真理的敌人——试图引述神圣记忆的西里尔和普罗克鲁斯的某些段落，似乎这些文字是支持特奥多雷的，因此正适合对他们引用先知的话：\Quote{上主的道路是正直的，义人在其中行走；但恶人在其中却要跌倒。}这些人恶毒地接受圣教父们所妥善合时地写下的东西，并在自己的罪中寻找借口，引用了这些话。教父们并非显得要免除特奥多雷的绝罚，反而是为了那些为聂斯多略及其不敬之心辩护的人，而经济地使用了某些表达，以便将他们从这错误中拉回，引领他们达至成全，并教导他们谴责不敬之心的门徒聂斯多略，以及其导师特奥多雷。因此，在这些经济的话中，教父们显示了他们的意图：特奥多雷应当受绝罚，正如我们从神圣记忆的西里尔和普罗克鲁斯关于谴责特奥多雷及其不敬之心的著作中，在本卷中已充分证明的。这种经济在圣经中也有发现：显然，宗徒保禄在传教初期，对那些在犹太人中长大的人使用了这种经济，并给弟茂德行了割损，为通过这种经济和迁就，引领他们达至成全。但后来他禁止了割损，如此写信给迦拉达人：\Quote{注意，我保禄告诉你们，如果你们受割损，基督对你们就毫无益处。}但我们发现，异端者们常做的事，特奥多雷的辩护者也做了。因为他们删去了圣教父们所写的某些内容，并掺入自己的一些虚假东西，试图通过神圣记忆的西里尔的一封信，好像来自教父们的证词，来免除上述不敬的特奥多雷的绝罚；而在这些段落中，当被删去的部分按其正确顺序读出时，真理便得到了证明，并且通过真实文本的比对，虚假被彻底揭露。但在所有这些事上，说这些虚妄话的人\Quote{信赖虚妄}，正如所记载的：\Quote{他们信赖虚妄，说虚妄的话；他们怀着痛苦，生出不义，编织蜘蛛网。}当我们如此考虑了特奥多雷及其不敬之心后，我们注意将戴奥多莱不敬地反对正统信仰、反对圣西里尔的\WorkTitle{十二章}和反对第一次厄弗所大公会议的一些内容，以及他为辩护不敬的特奥多雷和聂斯多略而写的一些东西，在我们卷中宣读并收录，以便读者满意；使众人知道这些人已被公正地驱逐和绝罚。第三，那封据说由依巴斯写给波斯人玛利斯信被提出审查，我们发现也应宣读。当宣读时，其不敬之心立刻为众人所明。而且，对上述\WorkTitle{三章}的谴责和绝罚是正确的，因为直到此时，对此仍有疑问。但由于这些不敬者的辩护者——特奥多雷和聂斯多略——正在设法以某种方式确认这些人和他们的不敬之心，并声称那封赞扬并辩护特奥多雷和聂斯多略及其不敬之心的不敬信函，被神圣的加采东大公会议所接纳，我们认为有必要表明，神圣会议免除了该信函中包含的不敬之心，以便澄清：说这些话的人并非得到该神圣会议的赞成，而是通过其名义来确认自己的不敬之心。而且在卷中已表明，依巴斯先前因这封信中含有的不敬之心而被指控：起初由神圣记忆的君士坦丁堡主教普罗克鲁斯，后来由虔诚记忆的德奥多西和弗拉维亚诺（他继普罗克鲁斯之后被祝圣为主教），他们将此事委托给提洛主教佛提乌和贝鲁特城主教欧斯塔基进行审查。后来，同一依巴斯被定罪，并被逐出主教职。既然如此，怎能有人胆敢说那封不敬信函被神圣的加采东大公会议接纳，并且神圣的加采东大公会议完全同意此信呢？然而，为使那些如此诽谤神圣加采东大公会议的人不再有机会这样做，我们命令宣读神圣会议——即第一次厄弗所会议和加采东会议——关于有福记忆的西里尔和虔诚记忆的旧罗马教宗良的书信的决定。并且，由于我们从这些得知，除已被证明与圣教父们的正统信仰一致者外，不应接纳任何人所写的任何东西，我们中断了进程，以便也宣读神圣加采东大公会议所颁布的信仰界定，从而将信函中的内容与这一定令进行对比。当做到这一点时，完全清楚地表明信函的内容与定令完全相反。\par

因为定令与由三百一十八位圣教父、一百五十位圣教父以及第一次厄弗所大公会议所颁布的一贯不变的信仰一致。但那封不敬的信函却相反，包含了异端者特奥多雷和聂斯多略的亵渎，并为他们辩护，称他们为导师，而称圣教父们为异端者。\par

我们将此向众人表明，我们绝无意图忽略第一次和第二次讨论的教父们，这些教父是戴奥多若和奈斯多利的追随者所引用来支持己方的；但在阅读了这些以及所有其他文件并审查其内容后，我们发现，前述的伊巴斯若不被强迫绝罚奈斯多利及其不敬的教导（这些教导在那封书信中受到辩护），是不能被接纳的。上述神圣大公会议的其他虔诚主教们，以及与那两位（某些人试图利用他们的讨论记录）也都如此行。\par

因为他们在戴奥多肋的案例中也遵守了此规则，并要求他绝罚他所被指控的那些事。因此，如果他们不允许以其他方式接纳伊巴斯，除非他谴责他书信中所包含的不敬，并签署大公会议所采纳的信德定义，他们又怎能试图声称这封不敬的书信被同一神圣大公会议所接纳呢？因为我们受到教导：\BibleQuote{正义与不法有什么相通？光明与黑暗有什么联系？基督与贝里雅耳有什么和谐？信者与不信者有什么相干？天主的殿与偶像有什么一致？}\BibleRef{格后 6:14-16}\par

在详细叙述了我们所做的一切之后，我们再次宣认：我们接纳四次神圣大公会议，即尼西大公会议、君士坦丁堡大公会议、第一次厄弗所大公会议和加采东大公会议，并且我们教导并继续教导它们关于唯一信德所界定的一切。凡不接纳这些的，我们视之为与圣而天主教会、宗徒的教会分离的人。此外，我们谴责并绝罚一切其他异端者，连同已被前述四次神圣大公会议以及圣而天主教会和宗徒的教会所谴责和绝罚的异端者，其中包括：莫普苏埃斯蒂亚的主教戴奥多若及其不敬的著作；以及戴奥多肋不敬地反对正确信德、反对圣济利禄的\WorkTitle{十二章}、反对第一次厄弗所大公会议所写的一切，以及他为戴奥多若和奈斯多利辩护所写的一切。此外，我们还绝罚那封据说是伊巴斯写给波斯人玛利斯的不敬书信，这封信否认天主圣言由天主之母、卒世童贞玛利亚降生成人，指责可纪念的济利禄（他教导真理）为异端者，并与他持有与阿波里纳相同的见解，且指责第一次厄弗所大公会议未经审查和调查就罢免了奈斯多利，称圣济利禄的\WorkTitle{十二章}为不敬的、违反正确信德，并辩护戴奥多若和奈斯多利及其不敬的教义和著作。因此，我们绝罚前述的三章，即：不敬的莫普苏埃斯蒂亚的戴奥多若及其可憎的著作；戴奥多肋不敬地所写的一切；以及据说是伊巴斯的不敬书信；以及它们的辩护者，那些已经或正在为它们辩护的人，或者胆敢声称它们是正确的人，以及那些已经或企图以圣教父或神圣加采东大公会议的名义为它们的不敬辩护的人。既然这一切都已准确裁定，我们谨记有关圣教会的许诺，以及那位说地狱之门不能战胜她（即异端者致命的舌头）的话；也记起\BibleBook{欧瑟亚}{书}中关于她的预言：\BibleQuote{我以忠实聘娶你，使你认识主。}\BibleRef{欧 2:20}我们将那些坚持不敬直至死亡的异端者放肆的舌头及其最不敬的著作，与谎言之父魔鬼一同归入行列，并对他们说：\BibleQuote{请看，你们点燃了火，使火焰旺盛；你们要在你们所燃的火光中行走，在你们所点燃的火焰中行走。}\BibleRef{依 50:11}但我们，既然受命以正确的教义劝勉民众，并对耶路撒冷（即天主的教会）的心说话，我们正确地急忙播种正义，收割生命的果实；并为自己从圣经和教父的教导中点燃知识之光，我们有必要以若干条文涵盖真理的宣告以及对异端者及其邪恶的谴责。\par

% structure: p0039 source=h2 target=section reason=relative-heading-rank; base=h2

\section{大公会议条文}

% structure: p0040 source=h3 target=subsection reason=relative-heading-rank; base=h2

\subsection{1}

若有人不宣认圣父、圣子、圣神的本性或本质、能力及权能是同一的；不宣认同体的天主圣三、一个天主性在三个位格中受钦崇，此人应受绝罚。因为只有一个天主，就是万物所从出的父；一个主耶稣基督，万物藉着他而存在；一个圣神，万物在他内存在。\par

% structure: p0042 source=h3 target=subsection reason=relative-heading-rank; base=h2

\subsection{2}

若有人不宣认天主的圣言有两次诞生：一次是从永远、由圣父所生，无始无终、无形无体；另一次是在末世，从天降下，由圣洁荣福的天主之母、卒世童贞玛利亚取得血肉，并由她所生，此人应受绝罚。\par

% structure: p0044 source=h3 target=subsection reason=relative-heading-rank; base=h2

\subsection{3}

若有人说行奇事的天主圣言是一位，而受苦的基督是另一位；或者说天主圣言与由女人所生的基督是分开的，或在他内如同一个人在另一个人内，而不是说我们的主耶稣基督，即降生成人、成为人的天主圣言，是同一的，他的奇迹和他按自己旨意在肉身内所受的苦难都属于同一者，此人应受绝罚。\par

% structure: p0046 source=h3 target=subsection reason=relative-heading-rank; base=h2

\subsection{4}

若有人声称天主圣言与人的结合仅是按照恩宠或动能、或尊严、或尊荣的平等、或权能、或关系、或效果、或权力、或按善意，即天主圣言悦纳一个人，也就是说，祂为了那人本身的缘故而爱他，如荒谬的德奥多若所言，或（若有人声称这种结合仅存在于）名称的相似上，如聂斯多略派所理解的，他们也称天主圣言为耶稣和基督，甚至将基督和圣子的名称加给那人，如此清楚地说有两个位格，却狡猾地只是在尊荣、尊严或敬拜方面指称一个位格和一位基督；若有人不承认如圣教父所教导的：天主圣言与有理性灵魂的肉身结合，并且这种结合是合成的和位格性的，因此只有一个位格，即我们的主耶稣基督，天主圣三中的一位：此人应受绝罚。正如事实所表明的，\Quote{联合}（\OriginalTerm{联合}{\GreekText{τῆς} \GreekText{ἑνώσεως}}）一词有多种含义，阿波利纳里和欧迪克的党羽声称这些性体\Emph{inter se}混淆，并主张一种由两者混合而产生的结合。另一方面，德奥多若和聂斯多略的追随者喜欢分裂性体，只教导一种相对性的结合。同时，天主的圣教会同等谴责这两种异端的亵渎，承认天主圣言与肉身的结合是合成性的，也就是说，位格性的。因为在基督的奥迹中，合成的结合不仅毫不混淆地保存了结合的性体，也不允许分离。\par

% structure: p0048 source=h3 target=subsection reason=relative-heading-rank; base=h2

\subsection{\Emph{5}}

若有人如此理解\Quote{我们的主耶稣基督只有一个位格}的说法，即认为这是许多本体的结合，并试图由此在基督的奥迹中引入两个本体或两个位格，且在引入两个位格后，仅按尊荣、尊敬或敬拜谈论一个位格，正如德奥多若和聂斯多略疯狂所写的那样；若有人诬蔑神圣的加采东大公会议，声称它以这种亵渎的意义使用了此说法（一个本体），而他不承认天主圣言是与肉身位格性地结合，并因此只有一个本体或一个位格，且神圣的加采东大公会议在这种意义上宣认了我们的主耶稣基督的一个位格：此人应受绝罚。因为天主圣三中的一位，即天主圣言，成了人，所以天主圣三并未因增加另一个位格或本体而扩大。\par

% structure: p0050 source=h3 target=subsection reason=relative-heading-rank; base=h2

\subsection{\Emph{6}}

若有人不以真义而仅以假义称圣、荣耀、永贞的玛利亚为天主之母，或仅在相对意义上如此称呼她，相信她只生下了一个普通人，而天主圣言并非由她降生成人，而是天主圣言由于与（由她所生的）那人结合而有了降生成人的结果；若他诬蔑神圣的加采东大公会议，好像它按照德奥多若的亵渎意义宣称了童贞女为天主之母；或若有人称她为人的母亲（\OriginalTerm{人的母亲}{\GreekText{ἀνθρωποτόκον}}）或基督的母亲（\OriginalTerm{基督的母亲}{\GreekText{Χριστοτόκον}}），好像基督不是天主，而他不确切真实地承认她是天主之母，因为那位在万世之前由圣父所生的天主圣言，在末世由她取得肉身并诞生，且若有人不承认神圣的加采东大公会议在这种意义上承认了她为天主之母：此人应受绝罚。\par

% structure: p0052 source=h3 target=subsection reason=relative-heading-rank; base=h2

\subsection{\Emph{7}}

若有人使用\Quote{在两个性体之内}的说法，却不承认我们的主耶稣基督一位在神性和人性中被启示出来，即以此说法指明那不可言喻的结合毫不混淆地构成的性体之间的区别，在此结合中，圣言的性体未变为肉身的性体，肉身的性体也未变为圣言的性体，因为各保留其本性，结合是位格性的；但若他按分裂各部分的意义来理解关于基督奥迹的说法，或在承认只有主耶稣（天主圣言成了人）内的两个性体时，不满足于理论性地把握那构成祂的性体之间的区别（此区别并未因两者间的结合而消失，因为一个由两者组成，两者在一个之内），反而利用数字（二）来分裂性体或使它们成为名副其实的位格：此人应受绝罚。\par

% structure: p0054 source=h3 target=subsection reason=relative-heading-rank; base=h2

\subsection{\Emph{8}}

若有人使用\Quote{出自两个性体}的说法，承认神性与人性结合了，或使用\Quote{天主圣言的一个性体成了肉身}的说法，而不像圣教父所教导的那样理解这些说法，即从神性和人性中有一个位格性的结合，由此形成一位基督；而是从这些说法试图引入一个由基督的神性和人性混合而成的性体或本体：此人应受绝罚。因为我们教导独生圣言是位格性地（与人性）结合，并非说性体彼此混淆，而是每个性体保留其所是，我们理解圣言与肉身结合了。因此，只有一位基督，同时是天主和人，按神性与圣父同性同体，按人性与我们同性同体。因此，那些分裂或分开基督神性救世计划奥迹的人，或在此奥迹中引入混乱的人，同样被天主的教会谴责和绝罚。\par

% structure: p0056 source=h3 target=subsection reason=relative-heading-rank; base=h2

\subsection{\Emph{9}}

若有人接受\Quote{基督应在他两个本性内受朝拜}的说法，其用意是引入两种朝拜——一种特别针对天主圣言，另一种则属于那人；或者，若有人为摆脱肉体（即基督的人性），或将神性与人性混合，妄言那结合（的本性）只有一个本性或本质（\OriginalTerm{本性即本质}{\GreekText{φύσιν} \GreekText{ἤγουν} \GreekText{οὐσίαν}}），并如此朝拜基督，而不以单一的朝拜敬拜降生成人的天主圣言连同其肉体，正如圣教会从起初所教导的：此人当受绝罚。\par

% structure: p0058 source=h3 target=subsection reason=relative-heading-rank; base=h2

\subsection{10}

若有人不承认在肉体中被钉死的主耶稣基督是真天主，是荣耀之主，并且是天主圣三中的一位：此人当受绝罚。\par

% structure: p0060 source=h3 target=subsection reason=relative-heading-rank; base=h2

\subsection{11}

若有人不谴责亚略、欧诺米、马其顿尼、阿波林、聂斯多略、欧迪奇及奥力振，连同他们不敬的著作，以及所有已被圣而公的宗徒的教会和前述四次神圣大公会议谴责与绝罚的其他异端者；并且（若有人）不同样谴责所有那些曾持有、或现在持有、或执迷不悟直至终末仍坚持与上述异端者相同意见的人：此人当受绝罚。\par

% structure: p0062 source=h3 target=subsection reason=relative-heading-rank; base=h2

\subsection{12}

若有人维护摩普绥提亚的不敬的德奥道若，他曾说：天主圣言是一个位格，而基督是另一个位格，为灵魂的痛苦和肉体的欲望所困扰，渐渐与卑下者分离，因善工的进步而变得更好，在生活方式上无可指摘，作为一个纯粹的人，因父、及子、及圣神之名受洗，并藉此洗礼获得圣神的恩宠，堪当承受子女的名分，并因天主圣言位格的缘故而被敬拜（正如人敬拜皇帝的肖像），且他在复活之后思想无可改变、全然无罪。而且，这同一个不敬的德奥道若还说：天主圣言与基督的结合，如同按宗徒的教导，男人与妻子之间的结合：\Quote{二人成为一体。} 同一个（德奥道若）竟胆敢在众多亵渎中说出：当主在复活后向门徒嘘气说：\Quote{你们领受圣神吧，}他并非真正赐予他们圣神，而只是以嘘气作为标记。他还说，多默在复活后触摸主的肋旁时所宣认的：\Quote{我的主！我的天主，}并非指着基督说的，而是多默因复活的奇迹而惊讶，如此感谢那使复活的天主。此外（更令人震惊的是），这同一个德奥道若在\WorkTitle{宗徒大事录}注释中将基督与柏拉图、摩尼、伊壁鸠鲁、玛尔基翁相比，并说：正如这些人各自发现自己的学说，并以自己的名字称呼门徒，称为柏拉图派、摩尼派、伊壁鸠鲁派、玛尔基翁派，同样基督发现了自己的学说，便将门徒称为基督徒。因此，若有人维护这最不敬的德奥道若及其不敬的著作，他在其中倾吐上述亵渎以及无数反对我们伟大的天主和救主耶稣基督的其他亵渎；若有人不谴责他及其不敬的著作，以及所有维护或捍卫他的人，或声称他的解释是正统的人，或为他及其不敬的著作写作的人，或持有相同意见的人，或曾持有并至死仍坚持此异端的人：此人当受绝罚。\par

% structure: p0064 source=h3 target=subsection reason=relative-heading-rank; base=h2

\subsection{13}

若有人维护德奥多肋的不敬著作，这些著作反对真信仰和第一次圣厄弗所大公会议，反对圣济利禄及其十二绝罚；并且（维护）他为维护不敬的德奥道若和聂斯多略以及持有与上述德奥道若和聂斯多略相同意见的人所写的；若有人接纳他们或其不敬之见，或将主张天主圣言位格结合的教会圣师称为不敬；若有人不谴责这些不敬的著作，以及那些曾持有或现在持有这些见解的人，和所有那些写作反对真信仰或反对圣济利禄及其十二章的人并且死于其不敬者：此人当受绝罚。\par

% structure: p0066 source=h3 target=subsection reason=relative-heading-rank; base=h2

\subsection{14}

若有人维护那封据说是依巴写给波斯人玛利斯的信函，在信中他否认由玛利亚——天主之母、终身童贞者——降生成人的天主圣言成了人，反而说由她所生的是一个纯粹的人，他称之为圣殿，仿佛天主圣言是一个位格，而那人又是另一个位格；在信中他还指责圣济利禄为异端者，而圣济利禄所教导的正是基督徒的正信，并指控他写出类似邪恶的阿波林的东西。此外，他还诋毁第一次圣厄弗所大公会议，断言会议未经甄别和审查就罢免了聂斯多略。上述不敬信函将真福济利禄的十二章称为与正信相悖的不敬之作，并维护德奥道若和聂斯多略及其不敬的教导和著作。因此，若有人维护前述信函而不谴责它以及那些维护它并声称它正确或部分正确的人；或者若有人维护那些已经或将要为它辩护的人，或维护为其中所含不敬之辞辩护的人，以及那些胆敢以圣教父或圣加采东大公会议的名义维护它或其中所含不敬之辞，并至死坚持这些罪行的人：此人当受绝罚。\par

% structure: p0068 source=h2 target=section reason=relative-heading-rank; base=h2

\section{针对奥力振的绝罚}

% structure: p0069 source=h3 target=subsection reason=relative-heading-rank; base=h2

\subsection{1}

若有人断言灵魂先存的神话，并断言由此产生的奇异复归论：此人当受绝罚。\par

% structure: p0071 source=h3 target=subsection reason=relative-heading-rank; base=h2

\subsection{2}

若有人说，一切理性受造物的受造（\OriginalTerm{受造}{\GreekText{τὴυ} \GreekText{παραγωγὴν}}）只包括无身体且完全非物质的理智者（\OriginalTerm{理智者}{\GreekText{νόας}}），既无数量亦无名称，以致因本质、能力与行动的同一，以及因他们与天主圣言的结合与认识而彼此统一；但他们不再渴望瞻仰天主，便随从自己的私欲而转向更坏的事物，且取了或粗或细的身体，并获得了名称，因为在天上品级之间，名称的差异正如身体的差异一样；从此有些成为并被称为革鲁宾，有些成为色辣芬、宰制者、异权者、德能者、上座者、天使，以及其他任何可能的天上品级：当受绝罚。\par

% structure: p0073 source=h3 target=subsection reason=relative-heading-rank; base=h2

\subsection{3}

若有人说，太阳、月亮和星辰也是理性受造物，它们之所以成为如今的样子，只是因为转向了邪恶：当受绝罚。\par

% structure: p0075 source=h3 target=subsection reason=relative-heading-rank; base=h2

\subsection{4}

若有人说，那些天主之爱变得冷淡的理性受造物，被隐藏在我们这样粗重的身体内，并被称为人；而那些达到了邪恶最低程度的，则分享了寒冷而黑暗的身体，并成为且被称为魔鬼和邪灵：当受绝罚。\par

% structure: p0077 source=h3 target=subsection reason=relative-heading-rank; base=h2

\subsection{5}

若有人说，来自天使或总领天使境界的是一种灵魂的（\OriginalTerm{灵魂的}{\GreekText{ψυχικὴν}}）状态，并且来自灵魂状态的又是魔鬼状态和人状态，并且来自人状态又可以再次成为天使和魔鬼，并且天上德能的每一等级都或是全部来自较低等级，或是全部来自较高等，或来自较高等和较低等：当受绝罚。\par

% structure: p0079 source=h3 target=subsection reason=relative-heading-rank; base=h2

\subsection{6}

若有人说，魔鬼有两类，一类包含人的灵魂，另一类包含堕落至此的较高神灵；并且在所有理性受造者的数量中，只有一个在爱和默观天主中保持不变，而这神灵成了基督，成了所有理性受造者的王，并创造了天上、地上和天地之间所有存在的身体；并且世界自身内含有比世界更古老的元素，这些元素自存，即：干、湿、热、冷，以及它所形成的理型（\OriginalTerm{理型}{\GreekText{ιδέαν}}）被如此形成；并且至圣且同质的天主圣三并未创造世界，而是被比世界更古老的造物主心智（\OriginalTerm{造物主心智}{\GreekText{Νοῦς} \GreekText{δημιρυργός}}）所创造，这心智将存在赋予世界：当受绝罚。\par

% structure: p0081 source=h3 target=subsection reason=relative-heading-rank; base=h2

\subsection{7}

若有人说，关于基督，据说祂以天主的形象出现，并在万世之前与天主圣言结合，并在末世贬抑自己以至人性，祂（按照他们的说法）怜悯了在统一同一性（祂自己是其中一部分）的神灵中出现的各种堕落，并为了恢复他们，经历了各种等级，有了不同身体和不同名字，成了万有中的万有，在天使中为天使，在异权者中为异权者，在德能者中为德能者，按不同等级理性受造者的形象穿上了相应等级的形状，最终取了和我们一样的血肉，为人而成了人；（若有人说这一切）而不承认是天主圣言贬抑自己并成了人：当受绝罚。\par

% structure: p0083 source=h3 target=subsection reason=relative-heading-rank; base=h2

\subsection{8}

若有人不承认天主圣言（与圣父和圣神同一实体，且成了血肉并成了人，是天主圣三之一）在任何意义上都是基督，而（断言）只是在一种不精确的意义上如此，并且是由于那被称为理智（\OriginalTerm{理智}{\GreekText{νοῦς}}）的空虚自己（\OriginalTerm{空虚自己}{\GreekText{κενώσαντα}}）的结果；若有人断言这与天主圣言联合的（\OriginalTerm{联合的}{\GreekText{συνημμένον}}）理智在真实意义上是基督，而圣言只是因为与理智的联合才被称为基督，并且反之，理智只是因为圣言才被称为天主：当受绝罚。\par

% structure: p0085 source=h3 target=subsection reason=relative-heading-rank; base=h2

\subsection{9}

若有人说，天主圣言降生成人，不是取了有灵魂（\OriginalTerm{灵魂}{\GreekText{ψυχὴ}}）的身体……\par

\SourceRecord{Translated by Henry Percival.；Nicene and Post-Nicene Fathers, Second Series；Vol. 14.；Edited by Philip Schaff and Henry Wace.；Buffalo, NY: Christian Literature Publishing Co.,；1900.；<http://www.newadvent.org/fathers/3812.htm>.}
